But, by pictures, the image of $F_{n+2}$ restricted to the
Temperley-Lieb subalgebra of  $N'\cap M_{n+1}$ is
$\dsize{\frac{1}{n+2}{2(n+1)\choose n+1}}$.
Hence on $TL^{\perp}$, \
dim$(F_1(TL^{\perp})=1$. So by Lemma 4.2.10,
dim$(TL^{\perp})\geq 2n+4$, a contradiction. \qed




We have obtained far more powerful results than the following by a
study of the representation theory of $AT(n,\d)$. These results
will be presented in a future paper of this series. We gave the
result here because it was announced some time ago. It is a version
of the ``triple point obstruction" of Haagerup and Ocneanu (see
[Ha]) but proved by a rather different method!

If we apply the argument we have just given when $\d \leq 2$
we obtain nontrivial but known results. The argument is very
simple so we present it.

{\bf Definition 4.2.12} The {\it critical depth} of a planar 
algebra $P$ will be the smallest $k$ for which there is an element
in $P_k$ which is not in the Temperley-Lieb subalgebra $TL_k$.

In the $C^*$- case, if $\d < 2$ ,the norm of the principal graph is less than $2$ so
as in [GHJ] it follows that the principal graph is an $A$,$D$ or
$E$ Coxeter graph with $*$ as far as possible from a vertex of 
valence $3$. In particular if $k$ is the critical depth, the 
dimension of the quotient $\frac {P_k}{V_k}$ is at most $1$ so the
rotation acts on it by multiplication by a $k$-th root of unity.
We will use the term ``chirality" for this root of unity in an
appropriate planar algebra.

The restrictions on the principal graph in the following theorem were first obtained by
Ocneanu.

{\bf Theorem 4.2.13} If $P$ is a $C^*$-planar algebra with $\d < 2$
then the principal graph can be neither $D_n$ with $n$ odd nor $E_7$.
If the principal graph is $D_{2n}$ the chirality is $-1$, if it
is $E_6$ the chirality is $e^{\pm 2\pi i/3}$ and if it is $E_8$ the 
chirality is $e^{\pm 2\pi i/5}$.


\bs
{\sc Proof.} If $k$ is the critical depth, by drawing diagrams one 
sees that the $\o$ in the $(2k+2)$x$(2k+2)$ matrix $\D_n(\o)$ is the chirality. Also if
$\d = z + z^{-1}$, we have  det$(\D_n(\o))= z^{2k+2} + z^{-(2k+2)} - \o - \o ^{-1}$.
If $\kappa$ is the Coxeter number of the principal graph we have 
$z = e^{\pm \pi i/{\kappa}}$. 
 
  On the other hand, by the argument of Corollary 4.2.11, the dimension
of $P_{k+1}$ would be too great if the determinant were non-zero. 
Thus we have, whatever the Coxeter graph may be,  

$$
{ e^{(2k+2)\pi i /\kappa} + e^{-(2k+2)\pi i /\kappa} = \o + \o ^{-1}}
$$

for a $k$-th root of unity $\o$.

The critical depth for $D_m$ would 
be $m-2$ so the left hand side of the equation is $-2$ so that $\o$ has to
be $-1$. But if $m$ is odd, $-1$ is not an $m-2$th root of unity so $D_m$ cannot
be a principal graph. If $m$ is even we conclude that the chirality is $-1$.

The critcal depth for $E_7$ is $4$ and the Coxeter number is 18. The above equation
clearly has no solution $\o$ which is a fourth root of unity.

For $E_6$ the critical depth is $3$ and the Coxeter number is $12$ so $\o = e^{\pm 2\pi i/3}$.
For $E_8$ the critical depth is $5$ and the Coxeter number is $30$ so $\o = e^{\pm 2\pi i/5}$. 

 
\qed
 

  The above analysis may also be carried out for $\d = 2$ where the Coxeter graphs are
replaced by the extended Coxeter graphs. In the $D$ case the presence of the Fuss Catalan 
algebra of example $2.3$ makes it appropriate to replace the notion of critical depth
by the first integer such that $P_k$ is bigger than the Fuss Catalan algebra. One obtains 
then that the chirality, together with the principal graph, is a complete invariant for
$C^*$-planar algebras with $\d = 2$ (see [EK] p. 586).
  
We end this section by giving more details of the planar structure
on $N\subset M$. In particular we give the subfactor
interpretations of duality, reduction, cabling and tensor product.
The free product for subfactors is less straightforward.

\bs\ni
{\bf Corollary 4.2.12} If $N\subset M$ is an extremal $II_1$
subfactor then, with the notation of $\S$3.2,
$\l_1(P^{N\subset M})=P^{M\subset M_1}$, as planar algebras.
(Note that $P^{M\subset M_1}_k=M'\cap M_{k}$ which is a subset of
$P^{N\subset M}_{k+1}$. We are saying that the identity map is an
isomorphism of planar algebras.)

\bs
{\sc Proof.} Equality of $P^{N\subset M}_k$ and
$\l_1(P^{N\subset M})$ as sets follows immediately from (ii) of
4.2.1. To show equality of the planar algebra structure we use the
uniqueness part of 4.2.1.

By definition of $\l_1(\Phi)$, for $x\in M'\cap M_k\subset N'\cap
M_k$,
\[
x=\Phi^{M\subset M_1}(\boxe{x})=
\Phi^{N\cap M}(\tangle{x}\ )=
\l_1(\Phi)(\boxe{x}\ ).
\]
Properties (i) and (iv) are straightforward as is the second
equation of (ii). So we only need to check the first equation of
(ii), i.e.,
\[
\l_1(\Phi^{N\subset M})( \vln\ \tanglel{x}\ )=\d E_{M_1}(x) 
\mbox{for}  x\in M'\cap M_k.
\]
But by definition
\[
\l_1(\Phi^{N\subset M})(\vln\ \tanglel{x}\ )
=\frac 1\d \Phi(\vln\  \vln\ \tanglel{x}\ ).
\]
If we define $g:\otimes^k_N M\to \otimes^k_N M$ and
$f:\otimes^k_N M\to\otimes^{k-2}_N M$ by
$g(y)=\sum_{b\in B} byb^*$ and
$f(x_1\otimes\dots\otimes x_k)=E(x_1)x_2\otimes\dots\otimes
x_{k-1} E(x_k)$, and if $\th(y)=x$ for $x\in M'\cap M_k$, by the
definition of $\Phi^{N\subset M}$ in Theorem 4.2.1,
\[
\Phi^{N\subset M}(\vln\ \vln\  \tanglel{x}\ )=\th
(\sum_{b\in B} b\otimes f(g(y))\otimes b^*).
\]
But since $\th$
is an $M-M$ bimodule map and $x$ commutes with $M$, \
$g(y)=\d^2y$. But by definition of $\th$ and 4.1.8,
\[
\th(\sum_{b\in B} b\otimes f(y)\otimes b^*)=
\d^{-2}\sum_{b\in B} bE_1 xE_1b^* \ .
\]
Since $\{bE_1\}$ is a basis for $M_1$ over $M$, we are done by
4.1.5. \qed

\bs Iterating, we see that $\l_n(P^{N\subset M})$ is the planar
algebra for the subfactor $M_{n-1}\subset M_n$.

For cabling we have the following

\bigskip
\ni{\bf Corollary 4.2.13} If $N\subset M$ is an extremal $II_1$
subfactor, the cabled planar algebra ${\cal C}_n(P^{N\subset M})$ of
$\S3.3$ is isomorphic to $P^{N\subset M_{n-1}}$.

\bigskip
{\sc Proof.} We will again use the uniqueness part of 4.2.1.
It follows from [PP2] that if we define $E^n_i$ to be
$$
(E_{ni}E_{ni-1}E_{ni-2}\dots E_{n(i-1)+1})
(E_{ni+1}E_{ni}\dots E_{n(i-1)+2})\dots
(E_{n(i+1)-1}E_{n(i+1)}\dots E_{ni})
$$
(a product of $n$ products of $E$'s with indices decreasing by
one), and $v^n_i$ with respect to $E^n_i$ as in 4.1.9, then the map
$x_1\otimes_N x_2\dots \otimes_N x_k\to x_1v^n_1 x_2v^n_2\dots
v^n_kx_k$ establishes, via the appropriate $\th$'s, a $*$-algebra
isomorphism between the $(k-1)$-th algebra in the tower for
$N\subset M_{n-1}$ and $M_{kn-1}$, hence between
$P_k^{N\subset M_{n-1}}$ and $P_{nk}^{N\subset M}=\cal C_n(P^{N\subset M})$.

If $\g$ is the inverse of this map, $\g\circ\cal C_n(\Phi)$ thus
defines a spherical planar algebra structure on
$P^{N\subset M_{n-1}}$. The labelling set is identified with
$P^{N\subset M_{n-1}}$ via $\g$, so
\[
\g\circ\Bbb C_n(\Phi^{N\subset M})(\tangle{x}\ )=
\Phi^{N\subset M})(\tangle{x})=x.
\]
Condition (i) of 4.1 for $\g\circ\Bbb C_n(\Phi)$ follows by
observing that\\ 
$\Bbb C_n(\Phi)(\boxj )=E^n_i$.
Condition (ii) follows from 4.2.12, and condition (iv) is clear.
\qed

\bs Reduction is a little more difficult to prove.

\bigskip
\ni{\bf  Corollary 4.2.14} Let $N\subset M$ be an extremal $II_1$
subfactor and $p$ a projection in $N'\cap M$. The planar algebra
$p(P^{N\subset M})p$ of $\S$3.3 is naturally isomorphic to the
planar algebra of the reduced subfactor
$pN\subset pMp$.

\bigskip
{\sc Proof.}  We first claim that the tautological map
$$
\a:\bigotimes^k_{pN} pMp\to\bigotimes^k_N M \ , \qquad
\a(\bigotimes^k_{i=1} x_i)=\bigotimes^k_{i=1} x_i
$$
is an injective homomorphism of $*$-algebras when both domain and
range of $\a$ are equipped with their algebra structures via the
respective maps $\th$ as in 4.1.10. To see this observe that the
conditional expectation $E_{pN}: pMp\to pN$ is just
$\frac{1}{\op{tr}(p)} E_N$. It follows that
$\a a_j=a_j\a$ and clearly $\a\mu_j=\mu_j\a$, so by 4.1.24,  $\a$
is a $*$-algebra homomorphism.

The next thing to show is that $\a$ takes $pN$-central vectors to
$\th^{-1}(p_k(N'\cap M_{k-1})p_k)$, with $p_k$ as in 3.3, using
the planar algebra structure on $N'\cap M_k$. Let
$\pi:\otimes^k_NM\to \otimes^k_N M$ be the map
$\pi(x_1\otimes x_2\dots\otimes x_k)=px_1p\otimes px_2p
\otimes \dots\otimes px_kp$, which is well defined since $p$
commutes with $N$. Then a diagram shows that, for
$x\in N'\cap M_{k-1}$,
$p_kxp_k=\th(\pi(\th(x)))$, and conversely, if
$\th^{-1}(x)$ is in the image of $\a$,
$\pi(\th^{-1}(x))=\th^{-1}(x)$ so $p_kxp_k=x$. Hence $\a$ induces
a $*$-algebra isomorphism between $P_k^{pN\subset pMp}$ and
$p_k(N'\cap M_{k-1})p_k$. To check that this map induces the
right planar algebra structure, we first observe that $\a$
commutes suitably with the maps $\eta$ and $\kappa$ of 4.1.19.
For $\eta$ we have $\a\circ\eta=\pi\circ\eta\circ\a$ by
definition. For $\kappa$, note that if we perform the basic
construction of [J1] on $L^2(M)$, \
$\frac{1}{\op{tr}(p)} pE_Np$ is the basic construction projection
for $pN\subset pMp$ on $L^2(pMp)$. Thus if $\{b\}$ is a basis for
$pMp$ over $pN$ we have
$\sum_b bE_N b^*={\op{tr}}(p)p$. So if
$$
x=x_1\otimes \dots\otimes x_k\in \bigotimes^k_{pN} pMp \ ,
$$
$$
\a(\kappa_j(x))=\frac{1}{\op{tr}(p)\d} \sum_b
x_1\otimes \dots\otimes x_j
b\otimes b^*\otimes x_{j+1}\otimes\dots\otimes x_k
$$
and
$$
\kappa_j(\a(x))=\frac 1\d \sum_c
x_1\otimes \dots\otimes x_j
c\otimes c^*\otimes x_{j+1}\otimes\dots\otimes x_k \ .
$$
Applying $\th$ we see that $\a \kappa_j=\kappa_j\a$.

With these two commutation results it is an easy matter to check
that $\a$ defines planar algebra isomorphism between
$P^{Np\subset pMp}$ and $p(P^{N\subset M})p$, using the uniqueness
part of 4.2.1 or otherwise.

\bigskip
\ni{\bf Corollary 4.2.15} If $N_1\subset M_1$ and
$N_2\subset M_2$ are extremal finite index subfactors, then
$P^{N_1\otimes N_2\subset M_1\otimes M_2}$ is naturally isomorphic
to the tensor product $P^{N_1\subset M_1}\otimes
P^{N_2\subset M_2}$ of $\S$3.4.

\bigskip
{\sc Proof.} We leave the details to the reader. \qed





\bs\bs\ni{\bf 4.3 \ Planar algebras give subfactors}

\smallskip
The following theorem relies heavily on a result of Popa [Po2].

\bigskip
\ni{\bf Theorem 4.3.1} Let $(P,\Phi)$ be a spherical $C^*$-planar
algebra with invariant $Z$ and trace ${\op{tr}}$. Then there is a
subfactor $N\subseteq M$ and isomorphisms $\Omega: N'\cap M_i\to
P_i$ with
\begin{verse}
(i) $\Omega$ is compatible with inclusions

(ii) ${\op{tr}}(\Omega(x))={\op{tr}}(x)$

(iii) $\Omega(M'\cap M_i)=P_{1,i}$
 (linear span of tangles with vertical first string)

(iv) $\Omega(e_i)=\frac{1}{\d}(\Phi(\bigtank ) )$

(v) $[M:N]=\d^2$. If $x\in N'\cap M_i$, choose $T\in\P$
with $\Phi(T)=\Omega(x)$.


(vi) $\Omega(E_{M'}(x))=\frac{1}{\d}\Phi (\boxk{T})$


(vii) $\Omega(E_{M_{i-1}}(x))=\frac{1}{\d}\Phi (\boxl{T})$
\end{verse}

\bigskip
{\sc Proof.} By theorem 3.1 of [Po2], the pair $N\subseteq M$
exists given a system $(A_{ij})$, $0 \leq i < j < \i$ of finite
dimensional $C^*$-algebras with $A_{i,j}\subset A_{k,\ell}$ if
$k\leq i$, $j\leq\ell$ and a faithful trace on $\bigcup^\i_{n=0}
A_{0n}$ satisfying 1.1.1, 1.1.2, 1.3.3$'$ and 2.1.1 of
[Po2]. We set $A_{i,j}=P_{i,j}$ (Definition 1.20). Then
$A_{ii}=\Bbb C$ since $Z$ is multiplicative and non-degenerate. The
conditions of %\linebreak% 
[Po2]  involve $e_i$'s and conditional
expectations
$E_{A_{ij}}$. We define the $e_i$'s in $P$ to be what we have called
$\Omega(e_i)$ (note our $e_i$ is
Popa's ``$e_{i+1}$"). The map $E_{A_{ij}}$ is defined by the
relation tr$(xE_{A_{ij}}(y))={\op{tr}}(xy)$ for $x$ in $A_{ij}$ and
$y$ arbitrary. Since $Z$ is an $S^2$ invariant one easily checks
that $E_{A_{ij}}$ is given by the element of $\cal A_{kj}(\emptyset)$ 
given in the figure below (for $x\in A_{0,k}$) 
\[
	\begin{picture}(0,0)%
\epsfig{file=ps/pic108.ps}%
\end{picture}%
\setlength{\unitlength}{0.00041700in}%
%
\begingroup\makeatletter\ifx\SetFigFont\undefined
% extract first six characters in \fmtname
\def\x#1#2#3#4#5#6#7\relax{\def\x{#1#2#3#4#5#6}}%
\expandafter\x\fmtname xxxxxx\relax \def\y{splain}%
\ifx\x\y   % LaTeX or SliTeX?
\gdef\SetFigFont#1#2#3{%
  \ifnum #1<17\tiny\else \ifnum #1<20\small\else
  \ifnum #1<24\normalsize\else \ifnum #1<29\large\else
  \ifnum #1<34\Large\else \ifnum #1<41\LARGE\else
     \huge\fi\fi\fi\fi\fi\fi
  \csname #3\endcsname}%
\else
\gdef\SetFigFont#1#2#3{\begingroup
  \count@#1\relax \ifnum 25<\count@\count@25\fi
  \def\x{\endgroup\@setsize\SetFigFont{#2pt}}%
  \expandafter\x
    \csname \romannumeral\the\count@ pt\expandafter\endcsname
    \csname @\romannumeral\the\count@ pt\endcsname
  \csname #3\endcsname}%
\fi
\fi\endgroup
\begin{picture}(7069,4369)(226,-4408)
\put(3901,-1486){\makebox(0,0)[lb]{\smash{\SetFigFont{12}{14.4}{rm}$\cdots$}}}
\put(4801,-1486){\makebox(0,0)[lb]{\smash{\SetFigFont{12}{14.4}{rm}$\cdots$}}}
\put(4801,-3736){\makebox(0,0)[lb]{\smash{\SetFigFont{12}{14.4}{rm}$\cdots$}}}
\put(3451,-3961){\makebox(0,0)[lb]{\smash{\SetFigFont{12}{14.4}{rm}$\cdots$}}}
\put(2851,-4336){\makebox(0,0)[lb]{\smash{\SetFigFont{12}{14.4}{rm}1}}}
\put(3376,-361){\makebox(0,0)[lb]{\smash{\SetFigFont{12}{14.4}{rm}$\cdots$}}}
\put(3151,-4336){\makebox(0,0)[lb]{\smash{\SetFigFont{12}{14.4}{rm}2}}}
\put(4126,-4336){\makebox(0,0)[lb]{\smash{\SetFigFont{12}{14.4}{rm}$i$}}}
\put(226,-2086){\makebox(0,0)[lb]{\smash{\SetFigFont{12}{14.4}{rm}$E_{A_{ij}}:\frac{1}{\delta^{i+k-j}}$}}}
\put(5851,-1486){\makebox(0,0)[lb]{\smash{\SetFigFont{12}{14.4}{rm}$\cdots$}}}
\put(5401,-4336){\makebox(0,0)[lb]{\smash{\SetFigFont{12}{14.4}{rm}$j$}}}
\put(4426,-4336){\makebox(0,0)[lb]{\smash{\SetFigFont{12}{14.4}{rm}$i+1$}}}
\end{picture}

\]
\ni Popa's  (1.1.1) and (1.1.2) and b)$'$ of (1.3.3) follow
immediately from pictures (note that the power of $\frac 1\d$ is
checked by applying $E_{A_{ij}}$ to 1). Condition a)$'$ of 1.3.3$'$
is dim $A_{ij}=
{\op{dim}} \ A_{i,j+1}e_j={\op{dim}} \ A_{i-1,j+1}$. But it is easy
from pictures that $E_{A_{ij}}$ defines a linear map from
$A_{i,j+1}e_j$ onto $A_{i,j}$, whose inverse is to embed in
$A_{i,j+1}$ and multiply on the right by $e_j$. Moreover th element
of $\cal A(\emptyset)$ (illustrated for $i\!=\!0$) in Figure 1
defines a linear isomorphism from $A_{i,j}$ to $A_{i+1,j+1}$ ---
the inverse is a similar picture.
\[
	\begin{picture}(0,0)%
\epsfig{file=ps/pic109.ps}%
\end{picture}%
\setlength{\unitlength}{0.00041700in}%
%
\begingroup\makeatletter\ifx\SetFigFont\undefined
% extract first six characters in \fmtname
\def\x#1#2#3#4#5#6#7\relax{\def\x{#1#2#3#4#5#6}}%
\expandafter\x\fmtname xxxxxx\relax \def\y{splain}%
\ifx\x\y   % LaTeX or SliTeX?
\gdef\SetFigFont#1#2#3{%
  \ifnum #1<17\tiny\else \ifnum #1<20\small\else
  \ifnum #1<24\normalsize\else \ifnum #1<29\large\else
  \ifnum #1<34\Large\else \ifnum #1<41\LARGE\else
     \huge\fi\fi\fi\fi\fi\fi
  \csname #3\endcsname}%
\else
\gdef\SetFigFont#1#2#3{\begingroup
  \count@#1\relax \ifnum 25<\count@\count@25\fi
  \def\x{\endgroup\@setsize\SetFigFont{#2pt}}%
  \expandafter\x
    \csname \romannumeral\the\count@ pt\expandafter\endcsname
    \csname @\romannumeral\the\count@ pt\endcsname
  \csname #3\endcsname}%
\fi
\fi\endgroup
\begin{picture}(4056,3637)(1231,-3983)
\put(3226,-961){\makebox(0,0)[lb]{\smash{\SetFigFont{12}{14.4}{rm}$\cdots$}}}
\put(3301,-3661){\makebox(0,0)[lb]{\smash{\SetFigFont{12}{14.4}{rm}$\cdots$}}}
\end{picture}

\]
\ni Finally  the commutation relations 2.1.1,
$[A_{ij},A_{k\ell}]=0$ for $i\leq j\leq k \leq\ell$ are trivial
since they are true in $\P$, involving non-overlapping strings.

It is standard theory for external subfactors that $E_{M'}$
restricted to $N'\cap M_i$ is $E_{A_{1,i}}$ and
$E_{M_{i-1}}=E_{A_{1,i-1}}$. So (vi) and (vii) are clear. \qed

\bigskip
\ni{\bf Corollary 4.3.2} If $(P,\Phi)$ is as before, the
Poincar\'e series $\sum^\i_{n=0}\dim (P_n)z^n$ has radius of
convergence $\geq\frac{1}{\d^2}$.

\bigskip
This result could be proved without the full strength of Theorem
4.3.1 (as pointed out by D.Bisch), using the principal graph and
the trace. The Poincar\'e series of planar subfactors enjoy many
special properties as we shall explore in future papers.

\newpage
\centerline{\bf References}

\begin{verse}
[A] V. Arnold, Remarks on the enumeration of plane curves.
{\it Amer. Math. Soc. Transl.} {\bf 173}(1996), 17--32.

[Bi] D.Bisch, Bimodules, higher relative commutants and the fusion algebra associated
to a subfactor,
{\it The Fields Institute for Research in Mathematical 
Sciences Communications Series} {\bf 13} (1997), 13--63.
 
[BW] J.Barrett and B.Westbury, ``Spherical Categories", hep-th/9310164.

[BiW] J.Birman and H.Wenzl, Braids, link polynomials and a new algebra.
{\it Trans. Amer. Math. Soc.} {\bf 163} (1989), 249--273.

[Ba] R. Baxter, {\it Exactly Solved Models in Statistical
Mechanics},   Academic Press, New York, 1982.

[Ban] T. Banica, Hopf algebras and subfactors associated to vertex models,
math.QA/9804016

[BH] D.Bisch and U.Haagerup, Composition of subfactors: new examples of infinite depth subfactors,
{\it Ann. scient. {\'E}c. Norm. Sup.} {\bf 29} (1996),329--383.

[BJ1] D.Bisch and V.Jones, ``Singly generated planar
algebras of small dimension",  preprint 1998.

[BJ2] D.Bisch and V.Jones, Algebras associated to
intermediate subfactors, {\it Invent. Math.} {\bf 128}, 89--157.
(1997).


[BJ3] D.Bisch and V.Jones, A note on free composition of
subfactors, Geometry and physics
 (Aarhus, 1995), 339--361 in {\it Lecture Notes in Pure and Appl.
Math.} {\bf 184}, Dekker, New York, 1997


[Co1] A.Connes, {\it Noncommutative Geometry}, Academic
Press (1994).

[dlH] P. de la Harpe,
Spin models for link polynomials, strongly regular graphs and Jaeger's Higman-Sims model.
{\it Pacific J. Math} {\bf 162}
(1994), 57--96.

[Dr] V.Drinfeld, Quantum groups, {\it Proceedings ICM
1986}, vol.~{\bf 1}, pp.798--820.


[EK] D.Evans and Y.Kawahigashi, {\it Quantum symmetries on operator algebras},
Oxford University Press (1998).

[F+] P.Freyd, D.Yetter, J.Hoste, W.Lickorish, K.Millett,
and A.Ocneanu, A new polynomial invariant of knots and links.
{\it Bull. AMS} {\bf 12} (1985), 183--190.


[FY] P.Freyd and D.Yetter,
 Braided compact closed categories with applications to low
dimensional topology, {\it Adv. in Math.} {\bf 77}
(1989), 156--182.


[FRS] K.~Fredenhagen, K-H.~Rehren and B.~Schroer,
Superselection sectors with braid group statistics and exchange
algebras, {\it Comm.~Math.~Phys.} {\bf 125} (1989), 201--226.



[GHJ] F.Goodman, P. de la Harpe and V.F.R.Jones, {\it
Coxeter Graphs and  Towers of Algebras}, Springer Verlag, MSRI
publications (1989).



[GL] J. Graham and G. Lehrer , ``The representation theory
of affine Temperley-Lieb algebras", preprint, University of Sydney
(1997).

[Gn] S. Gnerre, Free composition of paragroups,
Thesis, UC Berkeley, (1997).

[Ha] U.Haagerup,
Principal graphs of subfactors in the index range $4 < [M:N] <
3 + \sqrt{2}$
{\it Subfactors}, World Scientific, Singapore-New
Jersey-London-Hong Kong (1994), 1--39.



[JMN] F.Jaeger, M.Matsumoto and  K.Nomura, ``Bose-Mesner
algebras related with type  II matrices and spin models," RIMS
preprint (1995).

[Ja] F.JaegerStrongly regular graphs and spin models for the Kauffman polynomial, 
{\it Geom. Dedicata } {\bf 44} (1992), 23--52 


[J1] V.Jones,  Index for subfactors, {\it Invent.
Math.} {\bf 72}  (1983), 1--25.


[J2]  V.Jones, A polynomial invariant for knots via
von Neumann algebras.
       {\it Bulletin of the Amer. Math. Soc.} {\bf 12} (1985),
103--112.


[J3] V.Jones,  Hecke algebra representations of
braid groups and link polynomials, {\it Ann. Math.} {\bf 126}
(1987), 335--388.

[J4] V.Jones. On a certain value of the Kauffman polynomial,
{\it Comm. Math. Phys.} {\bf 125} (1989), 459-467

[J5] V.Jones Index for subrings of rings.
{\it Contemporary Math.} {/bf 43} (1985), 181--190.


[JS] V.F.R.Jones and V.Sunder, Introduction to
Subfactors, LMS lecture note
        series no.~{\bf 234} (1997) 162 pages.


[Ka1] L.Kauffman,
 State models and the Jones polynomial, {\it Topology} {\bf 26}
(1987), 395--407.

[Ka2] L.Kauffman,
An Invariant of regular isotopy, {\it Trans. Amer. Math. Soc.}
{\bf 318} (1990), 417--471.

[KS] Krishnan amd V.Sunder, On biunitary permutation matrices and some subfactors of index $9$, 
 {\it  Trans. Amer. Math. Soc} {\bf 348} (1996), 4691--4736.

[Ku] G.Kuperberg, Spiders for rank 2 Lie algebras, {\it
Commun. Math. Phys.} {\bf 180} (1996), 109--151, q-alg/9712003.

[La] Z. Landau, Thesis, UC Berkeley, 1998

[Lo] R.Longo, Index of subfactors and statistics of
quantum fields, I.
 {\it Comm. Math. Phys.} {\bf 126} (1989), 217--247.

[LM] W.Lickorish and K.Millett, Some evaluations of link 
polynomials.
{\it Comment. Math. Helv.} {\bf 61} (1986), 349--359.

[LR] R.Longo and J. Roberts, A Theory of Dimension,
{\it Journal of K-theory} {\bf 11} (1997), 103--159, funct-an/9604008.

[Ma] J.May, Definitions: operads, algebras and
modules, {\it Contemporary
              Mathematics } {\bf 202} (1997) 1--7.


[MvN] F.Murray and J.von Neumann, On rings of operators,
IV, {\it Annals Math.} {\bf 44} (1943)
      716--808.

[Mu] J.Murakami The Kauffman polynomial of links and representation theory,
{\it Osaka Journal of Mathematics} {\bf 24} (1987), 745--758.

[O1] A.Ocneanu,
 Quantum symmetry, differential geometry of finite graphs and
classification of subfactors,  University of Tokyo Seminary Notes
(notes recorded by Y. Kawahigashi).


[Pe] R.Penrose, Application of negative dimensional
tensors, in  {\it Combinatorial
    Mathematics and Its Applications}, Academic Press (1971)
221--244.


[Po1] S. Popa, Classification of subfactors and their
endomorphisms,  {\it CBMS Lecture Notes}, {\bf 86} (1995).



[Po2] S. Popa, An axiomatization of the lattice of higher
relative commutants of a subfactor, {\it Invent. Math.} {\bf120}
(1995), 427--445.


[PP] M.Pimsner and S.Popa, Entropy and index for
subfactors,  {\it Ann.
Sci. Ecole Norm. Sup.} {\bf19} (1986), 57--106.

[TL] H. Temperley and E.Lieb, Relations between the percolation....,
{\it Proc. Royal Soc. London} {\bf A322} (1971), 251--280.

[Tu] V.Turaev , {\it Quantum Invariants for knots and $3-$manifolds},
de Gruyter Berlin (1994).


[TV] V.Turaev and O.Viro, State sum invariants of
3-manifolds and quantum 6-j symbols. {\it
     Topology} {\bf 31} (1992) 86--902.

[Tut] W.Tutte, On dichromatic polynomials, 
{\it J. Combinatorial Theory} {\bf 2} (1967), 301--320.


[V] D.Voiculescu, Multiplication of certain non-commuting
random variables. {\it J. Operator Theory} {\bf 18} (1987),223--235.

[Wa] A.Wassermann, Operator algebras and conformal field theory, III:
Fusion of positive energy representations of LSU(N) using bounded
operators,
{\it Inventiones Math.} {\bf 133} (1998) 467--538, math.OA/9806031.

[We1]  H.~Wenzl,  Hecke algebras of type $A_n$ and
subfactors, {\it Invent.~Math.} {\bf 92} (1988), 249--383.

[We2]  H.~Wenzl,  Quantum groups and subfactors of type B, C
  and D, {\it Comm. Math. Phys.} {\bf 133} (1990), 383--432.

[We3]  H.~Wenzl, C* tensor categories from quantum groups,
{\ J. Amer. Math. Soc.} {\bf 11} (1998), 261--282. 

[X] F.Xu Standard $\lambda$-lattices from quantum groups,
{\it Inventiones Math.} {\bf 134} (1998),
\end{verse}

\enddocument


